\documentclass[10pt,a4paper]{amsart}
\usepackage[margin=19mm]{geometry}
\usepackage{amsmath,amssymb,mathtools}
\usepackage[scr=rsfs,cal=euler]{mathalfa}
\usepackage{xfrac}
\usepackage[unicode,hidelinks,bookmarks=false]{hyperref}
\newcommand{\ie}{i.e.\ }
\newcommand{\cf}{cf.\ }
\newcommand{\resp}{respectively\ }
\newcommand{\Subsection}{\subsection}
\providecommand{\nobreakdash}{\nobreak\hbox{-}}
\setlength{\emergencystretch}{2em}
\multlinegap0pt
\usepackage{bm}
\usepackage{bbm}
\usepackage{dsfont}
\usepackage{xspace}
\usepackage{cancel}
\usepackage{mathtools}
\usepackage{amsthm}
\makeatletter
\@ifundefined{thm}{\newtheorem{thm}{Theorem}}{}
\makeatother
\def\be{\begin{equation}}
\def\ee{\end{equation}}
\title[On local fields invariant under the action of topological de]{On local fields invariant under the action of topological defects}
\author{Anatoly Konechny, Vasileios Vergioglou}
\date{Source: arXiv:2505.04316v2 (CC BY 4.0)}
\begin{document}
\maketitle

\noindent\emph{Source and licence.} On local fields invariant under the action of topological defects. Version arXiv:2505.04316v2 (CC BY 4.0), distributed under the Creative Commons Attribution 4.0 International licence, as recorded in the arXiv licence metadata for this exact version and shown on the abstract page https://arxiv.org/abs/2505.04316v2. Full rights notice: licenses/arxiv-2505.04316-CC-BY-4.0.txt.\par\smallskip

\noindent\textit{Editorial scope.} Counted unit: the Thm whose source label is \texttt{su2thm} in the article (source0.tex lines 641--650, source label \texttt{su2thm}), delivered together with the article's own complete original proof of it (source0.tex lines 651--721, one \texttt{proof} environment of 71 lines). No mathematics is written, repaired or re-derived: statement and proof are the article's own text line for line. Required context retained uncounted: none. Every definition, display and label of those lines is retained unchanged. Omitted from this package and not redistributed: 1--640, 722--2588. Nothing inside a retained excerpt is omitted or altered: every retained body is delivered exactly as the article's lines are written, so no omission receipt of any kind accompanies this package. Citations: the retained excerpts carry no citation command at all, so no bibliography entry of the article is transcribed and no reference is redistributed. Reference adaptation: none was needed and none was made; every reference inside the delivered unit resolves inside it, so no unresolved reference or citation is produced. Printed slips kept literal: no printed word, symbol, hypothesis or number of the counted statement or of its proof was altered. Layout and class: the article's own class is \texttt{article} with packages including jheppub, lineno, fontenc, amsthm, amssymb, latexsym, amsfonts, pdfpages, float, tikz, gensymb, tikz-feynman, relsize, graphicx; the wrapper is the lane's pinned \texttt{amsart} editorial wrapper at 10pt on A4, which restates the theorem-like environments the retained text needs and adds the article's own macro definitions that the retained text needs (\texttt{\textbackslash be}, \texttt{\textbackslash ee}), with extra wrapper packages (bm, bbm, dsfont, xspace, cancel, mathtools). The delivered numbering is the unit's own. Assets: no retained span contains a figure environment, a graphics-inclusion command, a PDF-page inclusion, a table or a TikZ picture; no image file is copied into this package, the package declares no asset dependency and no image file is redistributed with it. Licence: version arXiv:2505.04316v2 of this article is distributed under the Creative Commons Attribution 4.0 International licence, as recorded in the arXiv licence metadata for this exact version and shown on the abstract page https://arxiv.org/abs/2505.04316v2. A full rights notice is delivered at licenses/arxiv-2505.04316-CC-BY-4.0.txt. Characters: every retained excerpt is pure ASCII, so the delivered engine, XeLaTeX with its default Latin Modern text font, reports no missing character.
\begin{thm}\label{su2thm}
   Let $a\in P_+^k\setminus \mathcal{J}_2$\footnote{If $(a)$ is any simple current then we have a trivial situation since $\mathcal{R}^k_{a,a}=\{(0)\}$.} and  $i,j\in P_+^k$ such that 
    \begin{equation}\label{con}
    i,j\in \mathcal{R}^k_{a,a}\,.
    \end{equation}
    Then if $i,j\notin \mathcal{J}_2$ the following holds
   \begin{equation}
       \operatorname{dim}\operatorname{Hom}\left((a)\times(a)\times(i)\times(j),(0)\right) >1 \,.
   \end{equation}
\end{thm}

\begin{proof}
Since the fusion product for  $\widehat{\mathfrak{su}}(2)_k$ representations is symmetric, we will assume without loss of generality that $i\geq j$. 
%Under the conditions of the theorem we have 
%    \begin{equation}
 %        \operatorname{dim}\operatorname{Hom}\left((a)\times(a)\times(i)\times(j),(0)\right)\neq 0\,.
  %  \end{equation}
%    Indeed, since every representation is self-conjugate we can rewrite condition \eqref{con} as $a\in \mathcal{R}^k_{i,a}$ and $a\in \mathcal{R}^k_{j,a}$. Then, consider the following two factors in the 4-product of interest
   % \begin{equation}\label{e1}
      %  (i)\times(a)=(i-a)\oplus\cdots\oplus(a)\oplus\cdots\oplus\operatorname{min}(i+a,2k-i-a)
   % \end{equation}
   % and
    %\begin{equation}\label{e2}
     %    (j)\times(a)=(j-a)\oplus\cdots\oplus(a)\oplus\cdots\oplus\operatorname{min}(j+a,2k-j-a).
    %\end{equation}
   % We see that there is a common representation $(a)$ in the r.h.s. of \eqref{e1} and \eqref{e2} which will lead to a copy of $(0)$ in the 4-product $(i)\times(a)\times(j)\times(a)$ coming from the term $(a)\times(a)$.
% Therefore we will prove the equivalent statement 
%\begin{equation}
%     \operatorname{dim}\operatorname{Hom}\left((a)\times(a)\times(i)\times(j),(0)\right) =1 \Longrightarrow \quad \text{at least one of }\, i,j\in \mathcal{J}_2\,.
%\end{equation}
By assumption we have  $0<i,j<k$ 
% arbitrary satisfying \eqref{con} 
and $a\notin\{0,k\}$ such that   %and consider the following two factors in the 4-product of interest
\begin{equation}\label{m}
        (a)\times(a)=(0)\oplus (2)\oplus\cdots(j)\oplus\cdots\oplus (i)\oplus\cdots \oplus \left(\operatorname{min}(2a,2k-2a)\right)
    \end{equation}
    where at least two representations must be present in the direct sum, 
    and
    \begin{equation}\label{ij}
        (i)\times(j)= (i-j)\oplus (i-j+2)\oplus\cdots \oplus \left(\operatorname{min}(i+j,2k-i-j)\right)
    \end{equation}
    The first expansion implies that  $i$ and $j$ are both even. Since $i\ge j$ and $j\ge 2$ the representations $(i-j)$ and $(i-j+2)$ must each appear in the right hand side of (\ref{m}). 
    While $(i-j)$ is also manifestly  present in the second expansion (\ref{ij}) to show that $(i-j+2)$ appears in (\ref{ij}) we need to check that under the assumptions at hand $i-j+2 \le \operatorname{min}(i+j,2k-i-j)$. Since $j\ge 2$ it follows that $i-j+2 \le i + j$. Furthermore, $i-j+2 \le 2k-i-j$ is equivalent to $i-j\le k-1$ that holds 
    due to the assumption $0<i,j<k$.   Hence, both representations $(i-j)$ and $(i-j+2)$ are present in expansion (\ref{m}) and in expansion (\ref{ij}) that implies 
    \be
     \operatorname{dim}\operatorname{Hom}\left((a)\times(a)\times(i)\times(j),(0)\right) >1 \, .
    \ee
   
    
%Since $a\notin\{0,k\}$ and $i-j\leq i$, the representations $(i-j),(i-j+2),\ldots$ appear on the r.h.s. of \eqref{m} with multiplicity one. The same representations also appear on the right hand side of \eqref{ij}. 

%Every representation is self-conjugate, therefore $\operatorname{dim}\operatorname{Hom}\left((a)\times(a)\times(i)\times(j),(0)\right)=1$ implies that there is only one common representation in the right-hand sides of \eqref{m} and \eqref{ij}. Since the representation $(i-j)$ is guaranteed to be in common, we deduce
%\begin{equation}\label{condi}
%    \operatorname{dim}\operatorname{Hom}\left((a)\times(a)\times(i)\times(j),(0)\right)=1\Longrightarrow i-j+2>\operatorname{min}\left(i+j,2k-i-j\right)
%\end{equation}
%To study the inequality
%\begin{equation}\label{ineq}
%    i-j+2>\operatorname{min}\left(i+j,2k-i-j\right)
%\end{equation}
%we consider both cases for the minimum as follows
%\begin{enumerate}
%        \item If 
 %       \begin{equation}\label{case1}
 %           i+j\leq 2k-i-j
  %      \end{equation}
  %  then \eqref{ineq} yields
  %      \begin{equation}
   %         i-j+2 > i+j \Longleftrightarrow j<1 \Longleftrightarrow j=0.
   %     \end{equation}
   %     Substituting $j=0$ back to \eqref{case1} we obtain $i\leq k$ which implies $i$ is arbitrary, so that the pair $j=0$ and $i$ arbitrary satisfies \eqref{ineq}.
   %     \item If
   %     \begin{equation}\label{case2}
    %        i+j\geq 2k-i-j
   %     \end{equation}
      %  then \eqref{ineq} yields
     %   \begin{equation}
     %       i-j+2 > 2k-i-j \Longleftrightarrow i>k-1 \Longleftrightarrow i=k.
     %   \end{equation}
     %   Substituting $i=k$ back to \eqref{case2} we obtain $j \geq 0$ which implies $j$ is arbitrary, so that the pair $i=k$ and $j$ arbitrary satisfies \eqref{ineq}.
   % \end{enumerate}
   % In total, \eqref{ineq} is satisfied if at least one  of $i,j$ are in $\{0,k\}$. Then from \eqref{condi} we conclude the proof.
\end{proof}

\end{document}
